% This file should be rename with the speaker's last name. (e.g : smith)

% Write the author names in the order you wish them to appear
% and underline the one who will give the talk.

% Only the speaker's email address is required.

\begin{abstract_online}{We have been told to use all the tools}{%
        \underline{Louis-Xavier Proulx}$^{1}$, Michel Beaudin$^{1}$, Anouk Bergeron-Brlek$^{1}$, Genevi\`eve Savard$^{1}$}
    
    {[louis-xavier.proulx@etsmtl.ca]}{$^1$  Service des enseignements g\'en\'eraux, \'Ecole de technologie sup\'erieure (\'ETS), Montr\'eal, Qu\'ebec, Canada}
    
In the mid nineties, when computer algebra systems and powerful symbolic calculators hit the market, the famous rule of three started to influence the teaching of some mathematics professors: ``\emph{every topic should be presented geometrically, numerically and algebraically}''~[1].  

In 1999, our engineering school decided to adopt Texas Instruments CAS technology.  During the following years,  many efforts were made in order to write lesson notes with new examples and exercises showing how this technology can be incorporated into the curriculum of some courses such as calculus [2] and differential equations [3]. These new ideas and approaches were presented and enhanced in dedicated education conferences, namely TIME, T$^{3}$ and ACA.

This talk will show concrete examples where graphical, numerical and symbolic tools can be used for solving problems in mathematics.

\textbf{References} \newline{}
% Model for references to books.
[1] \textsc{Deborah Hugues-Hallett, Andrew M. Geeason, et al.}, \emph{Calculus},  John Wiley \& Sons, 1994. \newline
[2] \textsc{Geneviève Savard, Robert Michaud, Andr\'e Bordeleau}, \emph{MAT145 - Calcul diff\'erentiel et int\'egral},  \small{\url{http://cours.etsmtl.ca/SEG/GSavard/mat145/exercices.html}}. \newline
[3] \textsc{Gilles Picard}, \emph{MAT265 - \'Equations diff\'erentielles}.  \small{\url{https://ena.etsmtl.ca/mod/page/view.php?id=101311}}. \newline
%  Model for references to articles

\end{abstract_online}
    